\begin{textsample}{POS Dim 7 – human – Score 4.00 – mg/mg\_ef\_ch\_geo\_7\_p5.txt}  \label{ex:f7_pos_004}
Entender a dimensão sociocultural e geopolítica da Eurásia na formação e constituição do Estado Moderno e nas disputas territoriais possibilita uma aprendizagem com ênfase no processo geo-histórico, ampliando e aprofundando as análises geopolíticas, por meio das situações geográficas que contextualizam os \textbf{temas} da geografia \textbf{regional}. \textbf{Espera} -se, assim, que o estudo da Geografia no Ensino Fundamental – Anos Finais possa contribuir para o delineamento do projeto de vida dos jovens estudantes, de modo que eles compreendam a produção social do espaço e a transformação do espaço em território usado.
Anseia- se, também, que entendam o papel do Estado-nação em um período histórico cuja inovação tecnológica é responsável por grandes transformações socioespaciais, acentuando ainda mais a necessidade de que possam conjecturar as alternativas de uso do território e as possibilidades de seus próprios projetos para o futuro. \textbf{Espera} -se, também, que, nesses estudos, sejam utilizadas diferentes representações cartográficas e linguagens para que os estudantes possam, por meio delas, entender o território, as territorialidades e o ordenamento territorial em diferentes escalas de análise.

% matched lemmas: esperar, regional, tema
\end{textsample}
